% ============================================================================
% backMatter.tex — conclusion generale, bibliographie/netographie
% ============================================================================

\chapter*{Conclusion generale}
\addcontentsline{toc}{chapter}{Conclusion generale}

Ce projet de fin d'annee nous a permis de concevoir et de realiser une preuve de
concept (PoC) complete d'une plateforme de streaming temps reel destinee a la
detection de fraude bancaire, en reponse aux limites du pipeline batch nocturne
actuellement en production chez HPS Group Morocco.

Sur le plan technique, nous avons mis en place un pipeline reposant exclusivement sur
des technologies open source : un cluster Apache Kafka gere par l'operateur Strimzi
pour le transport des evenements, quatre jobs Apache Flink (PyFlink) assurant
respectivement le decryptage, la validation selon neuf regles metier et ISO
bancaires, la normalisation et l'optimisation des transactions, une architecture de
stockage Medallion sur MinIO, et une persistance finale dans PostgreSQL via un
connecteur Kafka Connect JDBC Sink. L'ensemble de cette infrastructure est
provisionne de maniere reproductible via Terraform et supervise en temps reel par
Prometheus et Grafana.

Ce travail a egalement ete l'occasion de resoudre des difficultes techniques
concretes, notamment l'incompatibilite du connecteur JDBC Sink avec les messages JSON
schemaless --- resolue par l'introduction du composant \texttt{gold-flattener} --- ou
encore la necessite de figer certaines versions logicielles (Python 3.11) pour
garantir la stabilite des composants conteneurises.

\textbf{Limites du PoC.} Plusieurs limites doivent etre soulignees. Le stockage du
cluster Kafka et de PostgreSQL est actuellement de type \textit{ephemere}
(\texttt{emptyDir}), ce qui implique une perte de donnees en cas de redemarrage des
pods. Le cluster Flink fonctionne avec un seul JobManager, sans haute disponibilite.
Enfin, l'exporteur de metriques metier s'execute actuellement en dehors du cluster
Kubernetes, sur la machine hote.

\textbf{Perspectives.} Plusieurs ameliorations sont envisagees pour une mise en
production eventuelle :
\begin{itemize}
    \item migrer le stockage Kafka et PostgreSQL vers des \textit{Persistent Volume
    Claims} (PVC) afin de garantir la persistance des donnees au-dela du cycle de vie
    des pods,
    \item conteneuriser l'exporteur de metriques metier afin de l'integrer pleinement
    au cluster Kubernetes,
    \item deployer Apache Flink en haute disponibilite (plusieurs JobManagers avec
    election de leader),
    \item migrer l'infrastructure de Minikube vers un cluster Kubernetes manage tel
    qu'Azure Kubernetes Service (AKS), sans modification du code applicatif,
    \item explorer une resolution plus profonde de la limitation du schema JSON
    rencontree avec le connecteur JDBC Sink, par exemple via l'utilisation d'un
    registre de schemas (Schema Registry).
\end{itemize}

En definitive, ce projet demontre la viabilite technique d'un pipeline de detection de
fraude en temps reel construit entierement sur un socle open source, avec une latence
de traitement de l'ordre de quelques secondes contre plusieurs heures pour le systeme
batch existant.

\newpage

\chapter*{Bibliographie et Netographie}
\addcontentsline{toc}{chapter}{Bibliographie et Netographie}

\begin{thebibliography}{99}

\bibitem{kafka} Apache Software Foundation, \textit{Apache Kafka Documentation},
\url{https://kafka.apache.org/documentation/}.

\bibitem{flink} Apache Software Foundation, \textit{Apache Flink Documentation ---
PyFlink}, \url{https://nightlies.apache.org/flink/flink-docs-release-1.19/}.

\bibitem{strimzi} Strimzi, \textit{Strimzi Documentation --- Running Apache Kafka on
Kubernetes}, \url{https://strimzi.io/documentation/}.

\bibitem{minio} MinIO, \textit{MinIO Object Storage Documentation},
\url{https://min.io/docs/minio/linux/index.html}.

\bibitem{postgresql} PostgreSQL Global Development Group, \textit{PostgreSQL
Documentation}, \url{https://www.postgresql.org/docs/}.

\bibitem{debezium} Debezium, \textit{Debezium JDBC Connector Documentation},
\url{https://debezium.io/documentation/reference/stable/connectors/jdbc.html}.

\bibitem{prometheus} Prometheus, \textit{Prometheus Documentation},
\url{https://prometheus.io/docs/introduction/overview/}.

\bibitem{grafana} Grafana Labs, \textit{Grafana Documentation},
\url{https://grafana.com/docs/grafana/latest/}.

\bibitem{terraform} HashiCorp, \textit{Terraform Documentation},
\url{https://developer.hashicorp.com/terraform/docs}.

\bibitem{iso8583} International Organization for Standardization, \textit{ISO 8583 ---
Financial transaction card originated messages}.

\bibitem{iso4217} International Organization for Standardization, \textit{ISO 4217 ---
Currency codes}.

\bibitem{iso7812} International Organization for Standardization, \textit{ISO 7812 ---
Identification cards --- Identification of issuers}.

\bibitem{pcidss} PCI Security Standards Council, \textit{PCI Data Security Standard
(PCI-DSS)}, \url{https://www.pcisecuritystandards.org/}.

\bibitem{kubernetes} The Kubernetes Authors, \textit{Kubernetes Documentation},
\url{https://kubernetes.io/docs/home/}.

\end{thebibliography}
